\section{Implementation}
\label{sec:implementation}

\subsection{Module layout}

The implementation lives in
\code{zip-plus/src/code/binary\_add\_fft\_16/} and mirrors the
existing IPRS layout in \code{zip-plus/src/code/iprs/}:

\begin{itemize}
  \item \code{basis.rs}: a $\GF(2^{16})$ \code{Gf2\_16} type (a
    \code{u16} newtype) with addition (XOR), multiplication
    (carry-less multiply + reduce mod $f$), the half-trace solver
    (Gaussian elimination over $\F_2$) for $Y^2 + Y = \beta$, and
    Cantor basis generation.
  \item \code{params.rs}: a \code{Config16} trait declaring the
    extension degree and irreducible $f$, plus the precomputed
    Cantor basis, evaluation points, and subspace polynomials
    (stored sparsely as coefficient lists at the $\{2^j\}$
    positions).
  \item \code{ring\_ops.rs}: the lifted-ring arithmetic. Exports
    \code{reduce\_mod\_ftilde\_16}, the lifted-twiddle
    multiplication primitive, and the \code{FftRingElement16} trait.
  \item \code{radix8.rs}: the radix-8 evaluator ---
    \code{Radix8FftParams16}, which precomputes the per-meta-stage
    $8 \times 8$ Vandermonde tables, and \code{additive\_fft\_radix8\_16},
    the kernel.
  \item \code{ext\_field.rs}: the verifier-side field
    $F = \F_p[X]/\flift$ (\code{Gf2\_16Ext<P>}) and the
    $4$-bit-coefficient polynomial challenge type \code{Bit4Poly16}.
    See \autoref{sec:protocol-integration}.
  \item \code{packed\_cw.rs}: \code{PackedI64Poly16<BITS>}, a
    codeword/combined-row cell type with a bit-packed wire encoding
    (\autoref{sec:packed-cw}).
\end{itemize}

The top-level \code{binary\_add\_fft\_16.rs} exposes the public
\code{BinaryAddFft16Code} struct and its \code{LinearCode<Zt>} impl.

\subsection{The {\tt FftRingElement16} trait}
\label{sec:fft-ring-element}

The kernel is polymorphic over a \code{FftRingElement16} trait that
exposes the primitives the radix-8 matvec needs:

\begin{lstlisting}[language=Rust]
pub trait FftRingElement16: Sized + Clone + Send + Sync {
    fn fft_zero() -> Self;
    fn fft_add(&self, other: &Self) -> Self;
    fn fft_mul_lifted_gf16(&self, twiddle: Gf2_16) -> Self;

    // Batched-reduction path for the 8x8 matvec: accumulate the
    // eight lifted-twiddle products in an UNREDUCED buffer, then
    // reduce mod ftilde exactly once per output cell.
    type UnreducedAcc: Send + Sync + Clone;
    fn fft_unreduced_zero() -> Self::UnreducedAcc;
    fn fft_acc_mul_lifted_gf16(
        input: &Self, twiddle: Gf2_16, acc: &mut Self::UnreducedAcc,
    );
    fn fft_finalize_acc(acc: Self::UnreducedAcc) -> Self;
}
\end{lstlisting}

The \code{UnreducedAcc} associated type is the key to a fast radix-8
matvec: each $8 \times 8$ block accumulates all eight
lifted-twiddle products of an output cell into one length-$31$
unreduced buffer, and folds it back mod $\flift$ exactly once ---
rather than reducing after every one of the eight multiplications.

Implementations exist for \code{DensePolynomial<Int<N>, 16>} (any~$N$,
via crypto-bigint multi-limb arithmetic), the native
\code{DensePolynomial<i64, 16>} and \code{DensePolynomial<i128, 16>},
\code{Gf2\_16} itself (the unlifted reference), and
\code{PackedI64Poly16<BITS>} (which delegates to its inner
\code{DensePolynomial<i64, 16>}). The trait abstraction lets the
\code{ZipTypes} definition site choose the coefficient representation
without touching the kernel.

\subsection{The packed codeword cell}
\label{sec:packed-cw}

A codeword cell is a degree-$16$ polynomial. The shipped \code{Cw}
type is \code{PackedI64Poly16<BITS>}: in memory it is a
\code{DensePolynomial<i64, 16>} (so arithmetic uses native $64$-bit
ALU ops, with headroom for the FFT growth), but its \code{Transcribable}
(wire) encoding writes only \code{BITS} bits per coefficient instead of
the full $64$. The encoder packs $16$ coefficients of \code{BITS} bits
each into $\lceil 16 \cdot \code{BITS} / 8 \rceil$ bytes; the decoder
sign-extends each field back to \code{i64}.

This decouples the in-memory width (fixed at $64$ bits, for safe FFT
arithmetic) from the on-wire width (chosen per use). The shipped Zip$^+$
types use \code{BITS\_CW = 32} for codeword cells --- lossless given the
radix-8 codeword coefficients max near $2^{26}$, with $\sim 6$ bits of
headroom --- and \code{BITS\_CR = 8} for combined-row cells, whose
coefficients are bounded by $\code{batch} \cdot 15$ at $\code{num\_rows}
= 1$.

\subsection{The radix-8 kernel}
\label{sec:kernel}

\code{Radix8FftParams16::new(row\_len, codeword\_len)} asserts
$3 \mid \log_2(\code{codeword\_len})$ and precomputes the
per-meta-stage Vandermonde tables: a
\code{Vec<Vec<[[Gf2\_16; 8]; 8]>>} indexed by
\code{[meta\_stage][block\_idx]}, each entry an $8 \times 8$ matrix
of $\GF(2^{16})$ elements (later lifted at FFT time).

The kernel iterates the $m/3$ meta-stages from the top down; per
meta-stage it processes each block of $8$ in parallel; per block it
runs the $8 \times 8$ matvec using the batched-reduction path:

\begin{lstlisting}[language=Rust]
for j in 0..8 {
    let mut acc = T::fft_unreduced_zero();
    for i in 0..8 {
        let twiddle = matrix[block][j][i];
        if twiddle != Gf2_16::ZERO {
            T::fft_acc_mul_lifted_gf16(&input[i], twiddle, &mut acc);
        }
    }
    output[j] = T::fft_finalize_acc(acc);
}
\end{lstlisting}

\noindent Block-level parallelism (via \code{cfg\_chunks\_mut!} from
\code{zinc-utils}, which compiles to a rayon \code{par\_chunks\_mut}
under the \code{parallel} feature) is safe because each block is a
disjoint slice of the data vector. The zero-twiddle entries are
skipped --- in particular the $X_0 \equiv 1$ column contributes a
plain add.

\subsection{The {\tt LinearCode} impl}
\label{sec:linear-code-impl}

\code{BinaryAddFft16Code<Zt, C, REP, CHECK>} implements
\code{LinearCode<Zt>}:

\begin{lstlisting}[language=Rust]
impl<Zt, C, const REP: usize, const CHECK: bool> LinearCode<Zt>
    for BinaryAddFft16Code<Zt, C, REP, CHECK>
where
    Zt: ZipTypes,
    C: Config16,
    Zt::Cw: FftRingElement16 + FromRef<Zt::Eval>,
    Zt::CombR: FftRingElement16,
{
    const REPETITION_FACTOR: usize = REP;
    fn row_len(&self) -> usize { self.params.row_len }
    fn codeword_len(&self) -> usize { self.params.codeword_len }
    fn encode(&self, row: &[Zt::Eval]) -> Vec<Zt::Cw> {
        let mut data: Vec<Zt::Cw> = row.iter().map(Zt::Cw::from_ref).collect();
        data.resize_with(self.params.codeword_len, Zt::Cw::fft_zero);
        additive_fft_radix8_16(&mut data, &self.params);
        data
    }
    fn encode_wide(&self, row: &[Zt::CombR]) -> Vec<Zt::CombR> {
        let mut data: Vec<Zt::CombR> = row.to_vec();
        data.resize_with(self.params.codeword_len, Zt::CombR::fft_zero);
        additive_fft_radix8_16(&mut data, &self.params);
        data
    }
    fn encode_f<F>(&self, _row: &[F]) -> Vec<F> where ... {
        unimplemented!() // verifier uses encode_wide for proximity
    }
}
\end{lstlisting}

\code{encode} runs the radix-8 FFT on \code{Cw}-typed data;
\code{encode\_wide} runs the identical kernel on \code{CombR}-typed
data for the verifier's proximity check. \code{encode\_f} is part of
the \code{LinearCode} trait but unused by the Zip$^+$ protocol --- the
additive FFT over $\Rlift$ is not naturally expressible over an
arbitrary prime field --- so it is left \code{unimplemented!}.

\subsection{ZipTypes for the binary lane}
\label{sec:zip-types}

The shipped \code{ZipTypes} for this code sets \code{Comb} equal to
\code{CombR}; the prove/verify path then short-circuits the
alpha-sampling logic to $\alpha = [1]$ (because
$\code{Comb::DEGREE\_BOUND} = 0$ when $\code{Comb} = \code{CombR}$,
via a \code{Polynomial<Self>} blanket impl in
\code{poly/src/zero\_degree.rs}). The concrete shape used by the
bench and tests is

\begin{lstlisting}[language=Rust]
impl<...> ZipTypes for BenchAddFft16PolyChalPackedZipTypes<...> {
    const NUM_COLUMN_OPENINGS: usize = NUM_OPENINGS;
    type Eval  = BinaryPoly<16>;
    type Cw    = PackedI64Poly16<BITS_CW>;     // 32 bits/coef on wire
    type Fmod  = Uint<16>;                     // matches F::Modulus width
    type PrimeTest = MillerRabin;
    type Chal  = Bit4Poly16;                   // 4-bit-coef poly challenge
    type Pt    = Bit4Poly16;
    type CombR = PackedI64Poly16<BITS_CR>;     //  8 bits/coef on wire
    type Comb  = Self::CombR;                  // no alpha-projection
    type EvalDotChal     = ScalarProduct;
    type CombDotChal     = ScalarProduct;
    type ArrCombRDotChal = MBSInnerProduct;
}
\end{lstlisting}

\code{Eval = BinaryPoly<16>} is the natural shape of a $1\times$-folded
binary witness. \code{Cw} and \code{CombR} are both
\code{PackedI64Poly16} but with independent on-wire widths, since the
codeword coefficients are large (FFT-amplified) while the combined-row
coefficients stay small ($\le \code{batch}\cdot 15$ at
$\code{num\_rows} = 1$).
